% ============================================================================ 7 · Anexos
\section{Reproducibilidad: comandos y configuraciones}
\label{anx:repro}

Todo corre desde la raíz del repositorio con el entorno de \repo{requirements.txt} (más
\repo{requirements-dl.txt} para las redes). Los datos no se versionan: la entrega v2 va en
\repo{data/\allowbreak{}v2/\allowbreak{}raw/\allowbreak{}} y se apunta con \repo{FORD\_DATA\_DIR}. Sin red, \repo{WANDB\_MODE=disabled}. Cada número de
este informe sale de alguna de estas líneas. La GRU no es reproducible bit a bit entre plataformas ni cantidades de
hilos: una reproducción da números cercanos, no idénticos.

\subsection*{Datos, universo y panel}
\begin{lstlisting}
export FORD_DATA_DIR=$PWD/data/v2 WANDB_MODE=disabled OMP_NUM_THREADS=1
# entrega 1 vs v2
python scripts/compare_deliveries.py --config configs/data/compare_deliveries.yaml
# universo + holdout
python scripts/make_test_split.py --config configs/data/test_split.yaml --force
# cache del EDA (solo dev)
python scripts/build_eda_cache.py --config configs/data/eda_cache_v2.yaml
# EDA dentro del mercado
python scripts/eda_v2.py --config configs/eda_v2.yaml
# panel agregado (53 features)
python scripts/build_dataset.py --config configs/data/panel_v2.yaml
# + motor y serie
python scripts/build_dataset.py --config configs/data/panel_v2_estaticas.yaml
# folds 5 x 3, semilla 42
python scripts/make_splits.py --config configs/data/splits_panel_v2_r3.yaml
# secuencia 20 x 15
python scripts/build_seq_panel.py --config configs/data/panel_seq_trips_v2_estaticas.yaml
# chequeos del harness
python scripts/check_setup.py
\end{lstlisting}

\subsection*{El finalista en desarrollo}
\begin{lstlisting}
# semilla 42 (y _s1_r3, _s2_r3)
python scripts/train.py --config configs/exp_v2all_gru_trips_estaticas_r3.yaml
# ensamble por rango
python scripts/ensemble_rank.py --config configs/exp_v2all_seeds3_gru_trips_estaticas.yaml
# (a0), (a'), (b)
python scripts/audit_model.py --config configs/exp_v2all_gru_trips_estaticas_r3.yaml
# los 42 modelos, misma cuenta
python scripts/report_v2_models.py --config configs/report_v2_models.yaml
# anticipación en dev
python scripts/report_v2_leads.py --config configs/report_v2_leads.yaml
# sweep bayesiano (F10)
python scripts/sweep_gru.py --config configs/sweep_gru.yaml
# en qué se apoya
python scripts/explain_perm_seq.py --config configs/explain_gru_final.yaml
\end{lstlisting}

\subsection*{La medición en test}
\begin{lstlisting}
# el tiro preregistrado de F11
python scripts/eval_test.py --config configs/eval_test_f11.yaml --fit <yaml>
python scripts/eval_test.py --config configs/eval_test_f11.yaml --report
python scripts/eval_test.py --config configs/eval_test_f10_s42.yaml --fit configs/exp_v2all_gru_trips_estaticas_r3.yaml
# la GRU 42/1/2
python scripts/eval_test.py --config configs/eval_test_f10_s42.yaml --report
# todos los modelos
python scripts/eval_test.py --config configs/eval_test_modelos_v2.yaml --report
\end{lstlisting}

\subsection*{Este informe}
\begin{lstlisting}
# figuras y tablas generadas
python scripts/make_report_figures.py --config configs/report_figures.yaml
# PDF
cd docs/informe && tectonic informe.tex
\end{lstlisting}

\begin{table}[H]
\centering
\caption{Las configuraciones que definen el finalista.}
\small
\begin{tabularx}{\textwidth}{@{}>{\raggedright\arraybackslash}p{6.4cm}>{\raggedright\arraybackslash}X@{}}
\toprule
Archivo & Qué fija \\
\midrule
\repo{configs/\allowbreak{}data/\allowbreak{}raw\_sources.yaml} & Esquema de los crudos y correcciones de la entrega v2 \\
\repo{configs/\allowbreak{}data/\allowbreak{}vehicle\_dedupe.yaml} & Los 13 pares de clones \\
\repo{configs/\allowbreak{}data/\allowbreak{}test\_split.yaml} & Universo (ventana de producción) y sorteo del holdout \\
\repo{configs/\allowbreak{}data/\allowbreak{}panel\_v2.yaml}, \repo{panel\_v2\_estaticas.yaml} & W, G, H, Δ, referencia − 21 d, emparejado de
  sanos \\
\repo{configs/\allowbreak{}data/\allowbreak{}features\_v1.yaml} & Las 53 features, umbrales y topes físicos \\
\repo{configs/\allowbreak{}data/\allowbreak{}panel\_seq\_trips\_v2\_estaticas.yaml} & Los 15 canales y los 20 tramos de la secuencia \\
\repo{configs/\allowbreak{}exp\_v2all\_gru\_trips\_estaticas\{,\_s1,\_s2\}\_r3.yaml} & Arquitectura e hiperparámetros de la GRU, una
  por semilla \\
\repo{configs/\allowbreak{}exp\_v2all\_seeds3\_gru\_trips\_estaticas.yaml} & El ensamble por rango \\
\repo{configs/\allowbreak{}eval\_test\_f10\_s42.yaml}, \repo{eval\_test\_modelos\_v2.yaml} & La medición en test \\
\repo{configs/\allowbreak{}explain\_gru\_final.yaml} & La explicabilidad por permutación \\
\repo{configs/\allowbreak{}report\_figures.yaml} & Las figuras y tablas de este informe \\
\bottomrule
\end{tabularx}
\end{table}

% ----------------------------------------------------------------------------------------------
\section{Auditorías de leakage}
\label{anx:auditorias}

\textbf{Las tres defensas de construcción.} (1) El gap de blanking de 500 km sobre la referencia corrida 21 días;
(2) el split agrupado por vehículo, con los clones colapsados antes; (3) features solo hacia atrás y
preprocesamiento ajustado por fold. Además, el corte en el fin de extracción de los sanos impide que la sola
existencia de datos posteriores delate a un fallado.

\textbf{Las auditorías obligatorias} (\repo{scripts/\allowbreak{}audit\_model.py}), que toda corrida candidata tiene que pasar:
\begin{itemize}
  \item \textbf{(a0) nulo global}: se permutan las features entre todas las filas y se reentrena. El PR-AUC tiene
    que caer a la tasa base; si no cae, hay leakage.
  \item \textbf{\aprime{} aporte del cuándo}: se colapsa el puntaje de cada auto a su promedio, sin reentrenar. Lo
    que cae es lo que el modelo sabía del \emph{cuándo}.
  \item \textbf{(b) atajo de calendario}: se le dan al modelo las columnas \texttt{aux\_} de calendario (fechas de
    producción y venta, temperatura ambiente). Si el ROC sube más de 0,02, el modelo tiene un atajo disponible.
  \item \textbf{(a) permutación dentro del vehículo}: es \emph{informativa y no es un nulo}. Conserva qué autos
    fallan, así que sube en vez de bajar; no aprueba nada.
\end{itemize}

\begin{table}[H]
\centering
\caption{Auditorías de los modelos de la entrega v2 \tagdev{} (semilla 42 de cada secuencial). Tasa base de la
(a0): 0,058.}
\label{tab:auditorias}
\small
\begin{tabular}{@{}lrrrl@{}}
\toprule
Modelo & (a0) PR-AUC & \aprime{} & (b) calendario & Veredicto \\
\midrule
\textbf{GRU + viajes + estática (finalista)} & 0,061 & +0,001 & +0,004 & aprueba \\
GRU + viajes & 0,061 & +0,003 & \textbf{+0,025: marca} & aprueba \\
CNN-LSTM + viajes + estática & 0,057 & +0,009 & +0,018 & aprueba \\
CNN-LSTM + viajes & 0,060 & +0,026 & +0,001 & aprueba \\
CNN-LSTM de la tutora (señales + país) & 0,056 & +0,026 & \textbf{+0,028: marca} & aprueba \\
Survival stacking + motor y serie & 0,056 & +0,025 & +0,009 & aprueba \\
\bottomrule
\end{tabular}
\fuente{\repo{docs/\allowbreak{}memoria/\allowbreak{}f9-remedicion-completa-v2.md} (auditorías). Survival stacking sin estática: \aprime{}
+0,016 ± 0,004 (\repo{f9-remedicion-v2.md}).}
\end{table}

\textbf{Lectura.} Ningún modelo filtra: con las features permutadas, todos caen a la tasa base. La \aprime{} es
positiva en todos los modelos de la v2 pero chica: lo que separa es sobre todo \emph{qué auto}. El finalista no
marca el atajo de calendario; la misma GRU sin estática sí lo marcaba (+0,025), igual que la CNN-LSTM de la tutora
con solo señales.

\textbf{La auditoría que llevó a cambiar la métrica.} Con la primera entrega, el nulo global caía exacto a la tasa
base en tres modelos (pipeline limpio), pero el nulo dentro del vehículo quedaba \emph{por encima} de su propio
modelo (0,1928 contra 0,1653 del control): permutar dentro del auto conserva qué autos fallan y le saca el ruido de
qué ventana le tocó. Esa observación, junto con el techo de cohorte (Anexo~\ref{anx:prauc}), es la razón por la que
el finalista se elige por detección por auto, \aprime{} y estabilidad, y no por PR-AUC por fila
(\repo{decisiones.md}, cierre de F3).

\textbf{Una trampa de exposición que también auditamos.} Con la primera entrega, los fallados conservaban 18,25
cortes por bolsa y los sanos 9,00: el tamaño de la bolsa solo, como puntaje, daba lift 1,79×. Desde entonces toda
detección se compara contra un nulo que conserva el largo del historial (\repo{scripts/\allowbreak{}audit\_mil\_bagsize.py},
\repo{scripts/\allowbreak{}audit\_detection\_null.py}).

% ----------------------------------------------------------------------------------------------
\section{Todos los modelos en test}
\label{anx:modelos}

Cada modelo se entrenó con todo el desarrollo (semilla 42, salvo los ×3) y se midió con la misma cuenta: umbral
exacto por presupuesto de falsas alarmas, alerta de dos cortes seguidos, 103 autos (32 fallados). Se midieron
después del tiro preregistrado: describen, no eligen.

\begin{table}[H]
\centering
\caption{Detección (\%) en test por presupuesto de falsas alarmas \tagtest{}. “Prim.” es la media al 5–20\,\%.}
\label{tab:modelos-test}
\resizebox{\textwidth}{!}{% Generado por scripts/make_report_figures.py desde experiments/test-modelos-v2/curve.csv (no editar a mano).
\begin{tabular}{@{}lrrrrrrrrrl@{}}
\toprule
Modelo & 2\,\% & 5\,\% & 10\,\% & 15\,\% & 20\,\% & 25\,\% & 30\,\% & Prim. & AUC auto & dev (5 · 10 · 20\,\%) \\
\midrule
GRU ×3 (finalista) & 15,6 & 31,2 & 46,9 & 62,5 & 62,5 & 65,6 & 68,8 & \textbf{50,8} & 0,794 & 27 · 42 · 61 \\
CNN-LSTM ×3 & 3,1 & 15,6 & 56,2 & 56,2 & 68,8 & 68,8 & 68,8 & \textbf{49,2} & 0,790 & 22 · 41 · 62 \\
Survival stacking & 3,1 & 18,8 & 21,9 & 21,9 & 43,8 & 53,1 & 53,1 & \textbf{26,6} & 0,710 & 14 · 24 · 42 \\
LightGBM (53 features) & 0,0 & 15,6 & 25,0 & 28,1 & 34,4 & 53,1 & 56,2 & \textbf{25,8} & 0,725 & 18 · 28 · 42 \\
Logística L2 (53 features) & 0,0 & 6,2 & 12,5 & 31,2 & 34,4 & 46,9 & 53,1 & \textbf{21,1} & 0,679 & 15 · 23 · 37 \\
SS + efecto aleatorio (GPBoost) & 3,1 & 9,4 & 9,4 & 9,4 & 12,5 & 12,5 & 12,5 & \textbf{10,2} & 0,363 & 6 · 9 · 17 \\
\midrule
Azar, media / p95 (nulo de la GRU) & 3,4 / 9,4 & 7,0 / 15,6 & 13,7 / 25,0 & 18,6 / 31,2 & 24,4 / 40,6 & 28,6 / 43,8 & 33,9 / 50,0 & — & — & — \\
\bottomrule
\end{tabular}
}
\fuente{Detección y azar: \repo{experiments/\allowbreak{}test-modelos-v2/\allowbreak{}curve.csv} (tabla generada); AUC y dev:
\repo{docs/\allowbreak{}memoria/\allowbreak{}f11-test-resultado.md} §3.}
\end{table}

\textbf{Comparaciones pareadas} (bootstrap por vehículo, 2.000 réplicas, diferencia en puntos de primario):
\begin{itemize}
  \item GRU − celda mercado × motor: +27,3, IC95 [−0,8; 46,1], p(dif ≤ 0) = 0,028.
  \item Survival stacking − GRU: −24,2, IC95 [−43,0; −0,8] (no preregistrada).
  \item GRU 42/1/2 − GRU 101–103 (la misma configuración con otras semillas): +6,2, IC95 [−7,8; +18,0].
  \item GRU con etiqueta suave − GRU con etiqueta dura (el tiro preregistrado de F11, semillas 101–103 las dos):
    −13,3, IC95 [−25,0; +3,9].
\end{itemize}

\begin{table}[H]
\centering
\caption{El tiro preregistrado de F11 \tagtest{}: detección (\%).}
\label{tab:f11}
\scriptsize
\begin{tabular}{@{}lrrrrrrrr@{}}
\toprule
 & 2\,\% & 5\,\% & 10\,\% & 15\,\% & 20\,\% & 25\,\% & 30\,\% & Primario \\
\midrule
GRU con etiqueta suave (candidato) & 15,6 & 18,8 & 21,9 & 34,4 & 50,0 & 56,2 & 68,8 & 31,2 \\
GRU con etiqueta dura (referencia, 101–103) & 15,6 & 28,1 & 50,0 & 50,0 & 50,0 & 62,5 & 68,8 & 44,5 \\
Celda mercado × motor & 0,0 & 0,0 & 21,9 & 21,9 & 50,0 & 50,0 & 50,0 & 23,4 \\
\bottomrule
\end{tabular}
\fuente{\repo{docs/\allowbreak{}memoria/\allowbreak{}f11-test-resultado.md} §1; \repo{experiments/\allowbreak{}test-f11/\allowbreak{}}.}
\end{table}

\textbf{Notas.} La CNN-LSTM es inestable a presupuestos bajos (15,6\,\% al 5\,\%, 56,2\,\% al 10\,\%). GPBoost
invierte el orden: un auto nuevo entra sin su efecto aleatorio, y su AUC por auto ya estaba invertido en desarrollo
(0,34). La celda detecta 0\,\% por debajo del 10\,\% porque alerta celdas enteras.

% ----------------------------------------------------------------------------------------------
\section{Diccionario de features}
\label{anx:features}

\subsection*{Las 53 features agregadas (modelos tabulares y de supervivencia)}

Cada feature se calcula sobre la ventana (corte − 1.000 km, corte] del odómetro, solo con viajes y señales
anteriores al corte. Las \texttt{aux\_} quedan en el panel para auditar y no entran a ningún modelo. Agregadores:
\texttt{mean}/\texttt{median}/\texttt{p25}/\texttt{max} sobre los viajes de la ventana; \texttt{per\_1000km}, conteo
por 1.000 km recorridos; \texttt{per\_day}, por día calendario; \texttt{half\_*}, diferencia entre la segunda y la
primera mitad de la ventana (tendencia); \texttt{gap\_*}, distancia entre eventos; \texttt{km\_since\_last}, km desde
el último evento; \texttt{any}, si ocurrió alguna vez.

{\scriptsize
% Generado por scripts/make_report_figures.py desde configs/data/features_v1.yaml (no editar a mano).
% 53 features de modelo + 3 aux.
\begin{longtable}{@{}>{\raggedright\arraybackslash}p{5.4cm}l>{\raggedright\arraybackslash}p{4.6cm}>{\raggedright\arraybackslash}p{3.2cm}@{}}
\toprule
Feature & Fuente & Columna derivada & Agregador \\
\endhead
\midrule
\multicolumn{4}{@{}l}{\textbf{A · Térmica y trayectos cortos}} \\
\texttt{feat\_\allowbreak{}idle\_\allowbreak{}frac} & trips & \texttt{idle} & \texttt{mean} \\
\texttt{feat\_\allowbreak{}idle\_\allowbreak{}per\_\allowbreak{}1000km} & trips & \texttt{idle} & \texttt{per\_\allowbreak{}1000km} \\
\texttt{feat\_\allowbreak{}idle\_\allowbreak{}frac\_\allowbreak{}trend} & trips & \texttt{idle} & \texttt{half\_\allowbreak{}mean\_\allowbreak{}delta} \\
\texttt{feat\_\allowbreak{}short\_\allowbreak{}trip\_\allowbreak{}frac\_\allowbreak{}5km} & trips & \texttt{short\_\allowbreak{}trip\_\allowbreak{}moving} & \texttt{mean} \\
\texttt{feat\_\allowbreak{}trips\_\allowbreak{}below\_\allowbreak{}regime\_\allowbreak{}temp\_\allowbreak{}frac} & trips & \texttt{below\_\allowbreak{}regime} & \texttt{mean} \\
\texttt{feat\_\allowbreak{}below\_\allowbreak{}regime\_\allowbreak{}moving\_\allowbreak{}frac} & trips & \texttt{below\_\allowbreak{}regime\_\allowbreak{}moving} & \texttt{mean} \\
\texttt{feat\_\allowbreak{}below\_\allowbreak{}regime\_\allowbreak{}frac\_\allowbreak{}trend} & trips & \texttt{below\_\allowbreak{}regime} & \texttt{half\_\allowbreak{}mean\_\allowbreak{}delta} \\
\texttt{feat\_\allowbreak{}engine\_\allowbreak{}temp\_\allowbreak{}avg\_\allowbreak{}median} & trips & \texttt{EngineTemperatureAvg\_\allowbreak{}moving} & \texttt{median} \\
\texttt{feat\_\allowbreak{}engine\_\allowbreak{}temp\_\allowbreak{}max\_\allowbreak{}median} & trips & \texttt{EngineTemperatureMax\_\allowbreak{}moving} & \texttt{median} \\
\texttt{feat\_\allowbreak{}engine\_\allowbreak{}temp\_\allowbreak{}amplitude\_\allowbreak{}mean} & trips & \texttt{engine\_\allowbreak{}temp\_\allowbreak{}amplitude\_\allowbreak{}moving} & \texttt{mean} \\
\texttt{feat\_\allowbreak{}coolant\_\allowbreak{}temp\_\allowbreak{}end\_\allowbreak{}mean} & trips & \texttt{CoolantTemperatureEnd\_\allowbreak{}moving} & \texttt{mean} \\
\texttt{feat\_\allowbreak{}coolant\_\allowbreak{}temp\_\allowbreak{}start\_\allowbreak{}median} & trips & \texttt{CoolantTemperatureStart} & \texttt{median} \\
\texttt{feat\_\allowbreak{}cold\_\allowbreak{}start\_\allowbreak{}frac} & trips & \texttt{cold\_\allowbreak{}start} & \texttt{mean} \\
\texttt{feat\_\allowbreak{}chained\_\allowbreak{}trip\_\allowbreak{}frac} & trips & \texttt{chained} & \texttt{mean} \\
\texttt{feat\_\allowbreak{}trip\_\allowbreak{}distance\_\allowbreak{}median\_\allowbreak{}km} & trips & \texttt{trip\_\allowbreak{}km\_\allowbreak{}moving} & \texttt{median} \\
\texttt{feat\_\allowbreak{}trip\_\allowbreak{}distance\_\allowbreak{}p25\_\allowbreak{}km} & trips & \texttt{trip\_\allowbreak{}km\_\allowbreak{}moving} & \texttt{p25} \\
\texttt{feat\_\allowbreak{}trip\_\allowbreak{}duration\_\allowbreak{}median\_\allowbreak{}min} & trips & \texttt{trip\_\allowbreak{}duration\_\allowbreak{}min\_\allowbreak{}moving} & \texttt{median} \\
\midrule
\multicolumn{4}{@{}l}{\textbf{B · Regeneración (DPF)}} \\
\texttt{feat\_\allowbreak{}regenerations\_\allowbreak{}per\_\allowbreak{}1000km} & trips & \texttt{regen\_\allowbreak{}drop} & \texttt{per\_\allowbreak{}1000km} \\
\texttt{feat\_\allowbreak{}regenerations\_\allowbreak{}trend\_\allowbreak{}per\_\allowbreak{}1000km} & trips & \texttt{regen\_\allowbreak{}drop} & \texttt{half\_\allowbreak{}count\_\allowbreak{}delta\_\allowbreak{}per\_\allowbreak{}1000km} \\
\texttt{feat\_\allowbreak{}distance\_\allowbreak{}between\_\allowbreak{}regen\_\allowbreak{}mean\_\allowbreak{}km} & trips & \texttt{regen\_\allowbreak{}drop} & \texttt{gap\_\allowbreak{}mean\_\allowbreak{}km} \\
\texttt{feat\_\allowbreak{}distance\_\allowbreak{}between\_\allowbreak{}regen\_\allowbreak{}trend} & trips & \texttt{regen\_\allowbreak{}drop} & \texttt{gap\_\allowbreak{}trend\_\allowbreak{}km} \\
\texttt{feat\_\allowbreak{}km\_\allowbreak{}since\_\allowbreak{}last\_\allowbreak{}regen} & trips & \texttt{regen\_\allowbreak{}drop} & \texttt{km\_\allowbreak{}since\_\allowbreak{}last} \\
\texttt{feat\_\allowbreak{}regen\_\allowbreak{}residual\_\allowbreak{}mean} & trips & \texttt{regen\_\allowbreak{}residual} & \texttt{mean} \\
\texttt{feat\_\allowbreak{}regen\_\allowbreak{}start\_\allowbreak{}level\_\allowbreak{}mean} & trips & \texttt{regen\_\allowbreak{}start\_\allowbreak{}level} & \texttt{mean} \\
\texttt{feat\_\allowbreak{}dpf\_\allowbreak{}end\_\allowbreak{}mean} & trips & \texttt{AirRegenerationEnd} & \texttt{mean} \\
\texttt{feat\_\allowbreak{}dpf\_\allowbreak{}end\_\allowbreak{}max} & trips & \texttt{AirRegenerationEnd} & \texttt{max} \\
\texttt{feat\_\allowbreak{}dpf\_\allowbreak{}end\_\allowbreak{}slope\_\allowbreak{}per\_\allowbreak{}1000km} & trips & \texttt{AirRegenerationEnd} & \texttt{half\_\allowbreak{}slope\_\allowbreak{}per\_\allowbreak{}1000km} \\
\texttt{feat\_\allowbreak{}dpf\_\allowbreak{}positive\_\allowbreak{}delta\_\allowbreak{}frac} & trips & \texttt{dpf\_\allowbreak{}positive\_\allowbreak{}delta\_\allowbreak{}moving} & \texttt{mean} \\
\texttt{feat\_\allowbreak{}dpf\_\allowbreak{}saturated\_\allowbreak{}frac} & trips & \texttt{dpf\_\allowbreak{}saturated\_\allowbreak{}end} & \texttt{mean} \\
\texttt{feat\_\allowbreak{}filter\_\allowbreak{}abnormal\_\allowbreak{}end\_\allowbreak{}frac} & trips & \texttt{filter\_\allowbreak{}abnormal\_\allowbreak{}end} & \texttt{mean} \\
\texttt{feat\_\allowbreak{}filter\_\allowbreak{}cleaning\_\allowbreak{}auto\_\allowbreak{}end\_\allowbreak{}frac} & trips & \texttt{filter\_\allowbreak{}cleaning\_\allowbreak{}auto\_\allowbreak{}end} & \texttt{mean} \\
\texttt{feat\_\allowbreak{}manual\_\allowbreak{}regen\_\allowbreak{}ever} & trips & \texttt{filter\_\allowbreak{}manual} & \texttt{any} \\
\midrule
\multicolumn{4}{@{}l}{\textbf{C · Uso}} \\
\texttt{feat\_\allowbreak{}speed\_\allowbreak{}kmh\_\allowbreak{}mean} & trips & \texttt{speed\_\allowbreak{}kmh\_\allowbreak{}moving} & \texttt{mean} \\
\texttt{feat\_\allowbreak{}trips\_\allowbreak{}below\_\allowbreak{}30kmh\_\allowbreak{}frac} & trips & \texttt{urban\_\allowbreak{}moving} & \texttt{mean} \\
\texttt{feat\_\allowbreak{}moving\_\allowbreak{}trips\_\allowbreak{}per\_\allowbreak{}1000km} & trips & \texttt{moving} & \texttt{per\_\allowbreak{}1000km} \\
\texttt{feat\_\allowbreak{}km\_\allowbreak{}per\_\allowbreak{}day} & trips & \texttt{trip\_\allowbreak{}km} & \texttt{per\_\allowbreak{}day} \\
\texttt{feat\_\allowbreak{}trips\_\allowbreak{}per\_\allowbreak{}day} & trips & \texttt{moving} & \texttt{per\_\allowbreak{}day} \\
\texttt{feat\_\allowbreak{}hours\_\allowbreak{}between\_\allowbreak{}trips\_\allowbreak{}median} & trips & \texttt{soak\_\allowbreak{}min} & \texttt{median} \\
\texttt{aux\_\allowbreak{}air\_\allowbreak{}temp\_\allowbreak{}avg} & trips & \texttt{AirTemperatureAvg} & \texttt{mean} \\
\texttt{aux\_\allowbreak{}air\_\allowbreak{}temp\_\allowbreak{}min} & trips & \texttt{AirTemperatureMin} & \texttt{mean} \\
\midrule
\multicolumn{4}{@{}l}{\textbf{D · Severidad (mensajes, aceite, consumo)}} \\
\texttt{feat\_\allowbreak{}msg\_\allowbreak{}full\_\allowbreak{}per\_\allowbreak{}1000km} & signals & \texttt{msg\_\allowbreak{}full} & \texttt{per\_\allowbreak{}1000km} \\
\texttt{feat\_\allowbreak{}msg\_\allowbreak{}overloaded\_\allowbreak{}per\_\allowbreak{}1000km} & signals & \texttt{msg\_\allowbreak{}overloaded} & \texttt{per\_\allowbreak{}1000km} \\
\texttt{feat\_\allowbreak{}msg\_\allowbreak{}overloaded\_\allowbreak{}ever} & signals & \texttt{msg\_\allowbreak{}overloaded} & \texttt{any} \\
\texttt{feat\_\allowbreak{}msg\_\allowbreak{}over\_\allowbreak{}limit\_\allowbreak{}per\_\allowbreak{}1000km} & signals & \texttt{msg\_\allowbreak{}over\_\allowbreak{}limit} & \texttt{per\_\allowbreak{}1000km} \\
\texttt{feat\_\allowbreak{}msg\_\allowbreak{}over\_\allowbreak{}limit\_\allowbreak{}ever} & signals & \texttt{msg\_\allowbreak{}over\_\allowbreak{}limit} & \texttt{any} \\
\texttt{feat\_\allowbreak{}msg\_\allowbreak{}at\_\allowbreak{}limit\_\allowbreak{}ever} & signals & \texttt{msg\_\allowbreak{}at\_\allowbreak{}limit} & \texttt{any} \\
\texttt{feat\_\allowbreak{}msg\_\allowbreak{}cleaning\_\allowbreak{}auto\_\allowbreak{}per\_\allowbreak{}1000km} & signals & \texttt{msg\_\allowbreak{}cleaning\_\allowbreak{}auto} & \texttt{per\_\allowbreak{}1000km} \\
\texttt{feat\_\allowbreak{}msg\_\allowbreak{}cleaning\_\allowbreak{}auto\_\allowbreak{}trend\_\allowbreak{}per\_\allowbreak{}1000km} & signals & \texttt{msg\_\allowbreak{}cleaning\_\allowbreak{}auto} & \texttt{half\_\allowbreak{}count\_\allowbreak{}delta\_\allowbreak{}per\_\allowbreak{}1000km} \\
\texttt{feat\_\allowbreak{}msg\_\allowbreak{}abnormal\_\allowbreak{}frac} & signals & \texttt{msg\_\allowbreak{}abnormal} & \texttt{mean} \\
\texttt{feat\_\allowbreak{}msgs\_\allowbreak{}per\_\allowbreak{}1000km} & signals & \texttt{msg\_\allowbreak{}any} & \texttt{per\_\allowbreak{}1000km} \\
\texttt{feat\_\allowbreak{}oil\_\allowbreak{}life\_\allowbreak{}mean} & trips & \texttt{oil\_\allowbreak{}life} & \texttt{mean} \\
\texttt{feat\_\allowbreak{}oil\_\allowbreak{}life\_\allowbreak{}slope\_\allowbreak{}per\_\allowbreak{}1000km} & trips & \texttt{oil\_\allowbreak{}life} & \texttt{half\_\allowbreak{}slope\_\allowbreak{}per\_\allowbreak{}1000km} \\
\texttt{feat\_\allowbreak{}fuel\_\allowbreak{}pct\_\allowbreak{}per\_\allowbreak{}100km\_\allowbreak{}median} & trips & \texttt{fuel\_\allowbreak{}pct\_\allowbreak{}per\_\allowbreak{}100km} & \texttt{median} \\
\texttt{aux\_\allowbreak{}regen\_\allowbreak{}marker\_\allowbreak{}per\_\allowbreak{}1000km} & signals & \texttt{regen\_\allowbreak{}marker} & \texttt{per\_\allowbreak{}1000km} \\
\midrule
\multicolumn{4}{@{}l}{\textbf{Control de ventana (los emite \texttt{windows.py})}} \\
\texttt{feat\_n\_trips\_window} & trips & \texttt{*} & \texttt{count} \\
\texttt{feat\_window\_km\_covered} & trips & \texttt{OdometerTripEnd} & \texttt{rango} \\
\bottomrule
\end{longtable}

}
\fuente{\repo{configs/\allowbreak{}data/\allowbreak{}features\_v1.yaml} (tabla generada); derivadas en \repo{src/\allowbreak{}features/\allowbreak{}trips.py} y
\repo{src/\allowbreak{}features/\allowbreak{}signals.py}.}

\subsection*{Los 15 canales de la secuencia (la GRU)}

La ventana de 1.000 km se divide en 20 tramos de 50 km; un viaje cae en el tramo de su odómetro de fin. “zero”
rellena con cero un tramo sin datos (los conteos); “carry” arrastra el último valor dentro de la misma ventana,
siempre anterior al corte (un tramo sin viajes no es un tramo con ralentí cero).

\begin{center}
\scriptsize
% Generado por scripts/make_report_figures.py desde configs/data/panel_seq_trips_v2_estaticas.yaml (no editar a mano).
% ventana 1000 km en bins de 50 km = 20 tramos.
\begin{tabular}{@{}lllll@{}}
\toprule
Canal & Fuente & Columna & Agregado por tramo & Relleno \\
\midrule
\texttt{n\_records} & signals & \texttt{*} & count & zero · log1p \\
\texttt{observed} & signals & \texttt{*} & any & zero \\
\texttt{acum\_mean} & signals & \texttt{Acumulation} & mean & carry \\
\texttt{acum\_max} & signals & \texttt{Acumulation} & max & carry \\
\texttt{acum\_rise} & signals & \texttt{acum\_rise} & sum & zero \\
\texttt{regen\_drops} & signals & \texttt{acum\_regen\_drop} & sum & zero \\
\texttt{msg\_full\_frac} & signals & \texttt{msg\_full} & mean & zero \\
\texttt{msg\_severe\_frac} & signals & \texttt{msg\_severe} & mean & zero \\
\texttt{msg\_cleaning\_auto\_frac} & signals & \texttt{msg\_cleaning\_auto} & mean & zero \\
\texttt{n\_trips} & trips & \texttt{*} & count & zero · log1p \\
\texttt{idle\_frac} & trips & \texttt{idle} & mean & carry \\
\texttt{below\_regime\_frac} & trips & \texttt{below\_regime\_moving} & mean & carry \\
\texttt{speed\_kmh} & trips & \texttt{speed\_kmh\_moving} & mean & carry \\
\texttt{trip\_km} & trips & \texttt{trip\_km\_moving} & mean & carry \\
\texttt{trip\_duration\_min} & trips & \texttt{trip\_duration\_min\_moving} & mean & carry \\
\bottomrule
\end{tabular}

\end{center}
\fuente{\repo{configs/\allowbreak{}data/\allowbreak{}panel\_seq\_trips\_v2\_estaticas.yaml} (tabla generada). Más país, motor y serie en
one-hot como \texttt{static\_}.}

% ----------------------------------------------------------------------------------------------
\section{Las tablas de Ford: anexo contra archivos}
\label{anx:tablas}

\begin{table}[H]
\centering
\caption{\emph{Trip Summary}: las 40 columnas del anexo 7.1 y lo que llegó.}
\label{tab:trips-cols}
\scriptsize
\begin{tabularx}{\textwidth}{@{}>{\raggedright\arraybackslash}p{5.6cm}>{\raggedright\arraybackslash}p{2.6cm}>{\raggedright\arraybackslash}X@{}}
\toprule
Columna del anexo & Llegó & Uso \\
\midrule
\texttt{Vin} & como \texttt{VehicleCode} & Identificador anonimizado \\
\texttt{TripDatetimeStart/\allowbreak{}End} & sí & Duración, reposo entre viajes, emparejado por mes (dos formatos mezclados) \\
\texttt{OdometerTripStart/\allowbreak{}End} & sí & Eje de km, km por viaje, ventanas \\
\texttt{KilometerPerHour} & sí & Nulo en los viajes de 0 km; se recalcula como km/duración \\
\texttt{FuelLvlStartPc/\allowbreak{}EndPc} & sí & Consumo como índice relativo (supera 100) \\
\texttt{FuelLvlAutonomyStart} & no & — \\
\texttt{EngineOilLifePCStart/\allowbreak{}End} & sí & Pendiente del nivel en la ventana (no cambia dentro del viaje) \\
\texttt{EngineTemperatureMin/\allowbreak{}Max/\allowbreak{}Avg} & sí & Bajo régimen (máx. < 70 °C), temperaturas \\
\texttt{CoolantTemperatureStart/\allowbreak{}End} & sí & Arranque en frío, refrigerante al terminar \\
\texttt{DieselParticulateFilterStart/\allowbreak{}End} & no, por ese nombre & Llega como \texttt{AirRegenerationStart/\allowbreak{}End} (= \texttt{Acumulation}, 99,8\,\%) \\
\texttt{VehicleGPSLat/\allowbreak{}LongDataStart/\allowbreak{}End} (4) & no & Sin GPS no hay ubicación ni altura por viaje \\
\texttt{VehicleElevationRangeStart/\allowbreak{}End} & no & Proxy probado: altura de la ciudad de venta \\
\texttt{TirePressure\{LF,RF,LR,RR\}\{Start,End\}} (8) & no & — \\
\texttt{AirTemperatureStart/\allowbreak{}End/\allowbreak{}Min/\allowbreak{}Max/\allowbreak{}Avg} & sí & Solo auditoría: deriva con el mes y el lugar \\
\texttt{ManualRegenerationSootStart/\allowbreak{}End} & no & Los mensajes de limpieza manual marcan el evento (rarísimos) \\
\texttt{TripNumber} & sí & — \\
\emph{No está en el anexo}: \texttt{AirFilterStart/\allowbreak{}End} & sí & Nivel categórico del filtro (= \texttt{Message}, 99,97\,\%) \\
\bottomrule
\end{tabularx}
\fuente{\repo{docs/\allowbreak{}memoria/\allowbreak{}f2-diccionario-trips-incompleto.md}, \repo{f2-columnas-airfilter-airregeneration.md};
\repo{configs/\allowbreak{}data/\allowbreak{}raw\_sources.yaml}.}
\end{table}

\begin{table}[H]
\centering
\caption{\emph{DynamicInformation} y \emph{Vehicle General Information}.}
\label{tab:otras-cols}
\scriptsize
\begin{tabularx}{\textwidth}{@{}>{\raggedright\arraybackslash}p{4.6cm}>{\raggedright\arraybackslash}X@{}}
\toprule
Columna & Qué encontramos y cómo se usa \\
\midrule
\texttt{VehicleCode} & Clave de las seis tablas; 13 vehículos con dos códigos \\
\texttt{EventTimestamp} & Llega como \texttt{eventTimestamp}, dos formatos mezclados \\
\texttt{OdometerValue} & Nulos concentrados en pocos autos; la ventana se ubica por el odómetro de los viajes \\
\texttt{Acumulation} & Nivel del filtro, escala 0–100 en pasos de 5; la misma variable que \texttt{AirRegeneration} \\
\texttt{Message} & 9 niveles del filtro (“Normal Operation”, “Full”, “Overloaded”, “Over Limit”, limpiezas…); tasas
  por 1.000 km como severidad \\
\texttt{Regenerations} & El texto \texttt{“Regeneration”} marcando una fila; se corta el 25-05-2026: solo auditoría \\
\texttt{DistanceBetweenRegenerations} & Distancia al marcador anterior; hereda el corte: solo auditoría \\
\midrule
\texttt{VehCode} & Como \texttt{VehicleCode}; en v2 los fallados traen una fila por evento \\
\texttt{IdentificationDate} & Días desde producción del evento; nula en todos los sanos: es la etiqueta. En los
  fallados de v2 llega renombrada \\
\texttt{daysUntilSale} & Solo auditoría (cohorte) \\
\texttt{ProductionDay} & En el eje del calendario con origen común (IQR 0 días); en los fallados de v2, +538 días;
  solo auditoría \\
\texttt{Engine}, \texttt{ModelSeries} & Candidatas en v2; se reporta siempre dentro de mercado × motor \\
\texttt{SalesCountryCd} & Llega como \texttt{SalesCountry\_cd}; en v2 con código ISO \\
\emph{Nueva en v2}: \texttt{SalesCity} & Altura de la ciudad de venta (probada, no entra) \\
\bottomrule
\end{tabularx}
\fuente{\repo{f1-datos-reales.md}, \repo{f1-anclaje-temporal.md}, \repo{f9-entrega-v2.md}.}
\end{table}

% ----------------------------------------------------------------------------------------------
\section{Anexo técnico: el PR-AUC por fila y el techo de cohorte}
\label{anx:prauc}

Lo dejamos acá, y no en el cuerpo, porque en este panel no mide lo que Ford necesita: un PR-AUC por fila por
debajo del techo de cohorte no demuestra anticipación.

\begin{table}[H]
\centering
\caption{PR-AUC por fila en el dev v2 \tagdev{}, al lado del techo de cohorte.}
\label{tab:prauc}
\small
\begin{tabular}{@{}lrrr@{}}
\toprule
 & PR-AUC (lift) & ROC por fila & \aprime{} \\
\midrule
Tasa base & 0,057 (0,97×) & 0,489 & −0,002 \\
LightGBM (control) & 0,112 (1,92×) & 0,690 & −0,019 \\
Survival stacking & 0,125 (2,14×) & 0,709 & +0,016 \\
GRU + viajes + estática ×3 & 0,160 & — & +0,006 \\
\midrule
\textbf{Techo de cohorte} (“¿este auto falla?”, sin nada del cuándo) & \textbf{0,177 (3,03×)} & — & — \\
\bottomrule
\end{tabular}
\fuente{\repo{docs/\allowbreak{}memoria/\allowbreak{}f9-remedicion-v2.md}; GRU: \repo{f9-remedicion-completa-v2.md} (la \aprime{} del
ensamble es +0,006; la de la semilla 42 sola, +0,001).}
\end{table}

\textbf{Cómo se lee.} Todas las filas positivas del panel están dentro de los cortes de autos que fallan. Puntuar
cada fila con la cohorte del auto ya da 0,177; ningún modelo llega. Lo que un modelo agrega sobre ese techo se mide
con \aprime{} y con el PR-AUC restringido a los cortes de autos que fallan. Con la primera entrega el techo era
0,263 (2,10×), y la meta del plan inicial (lift de 1,6–2×) quedaba toda por debajo.
